\section{Introduction}

When we designed this study, we expected to find that supervised fine-tuning (SFT) fails on hard competitive programming benchmarks because the training dataset lacks correct solutions --- a plausible story backed by published performance numbers. We were wrong about the mechanism. The APPS dataset contains reference solutions for 85.32\% of competition-level problems. Yet a 7B code language model fine-tuned with SFT on those solutions achieves \textit{0.0 pass@1} on LiveCodeBench-Hard. This is not a dataset void. It is a generalization void.

This observation reframes a longstanding assumption in the reinforcement learning from execution feedback (RLEF) literature. Prior work argued that RLEF's advantage over SFT at hard difficulty stems from partial-success reward enabling gradient signal on examples where SFT has no correct training data. Our results suggest the mechanism is different: SFT has the training data but fails to generalize it to held-out hard evaluation problems (the precise mechanism remains under investigation; see \S6.2, L2). RLEF, by training on execution feedback from the model's own generated outputs, naturally operates at the model's actual capability frontier rather than an idealized reference distribution the model cannot reach.

The practical implications are significant. Practitioner decisions about whether and how to deploy RLEF --- particularly the choice of reward formulation --- rest on mechanistic assumptions that our controlled study allows us to examine directly. If the advantage is generalization-driven rather than data-coverage-driven, then the specific reward formulation may matter less than the execution feedback itself. This prediction is confirmed: fraction-of-tests and binary RLEF rewards converge to equivalent performance in our controlled experiment ($\Delta$=+0.0072, p=0.552). The key engineering investment is execution infrastructure, not reward design.

Despite the practical importance of RLEF for code generation, no prior controlled, reproducible, open-source study has directly compared RLEF against SFT across the full benchmark difficulty spectrum --- from HumanEval (easy) through LiveCodeBench-Hard (competitive programming). Existing work uses internal, non-reproducible models \cite{Gehring2024RLEF} or tests only on easy/medium benchmarks \cite{Le2022CodeRL,Majeed2023PPOCoder,Liu2023RLTF}. This reproducibility gap makes it impossible to verify difficulty-scaling claims or derive principled guidance for practitioners.

We fill this gap with a controlled study using publicly available components throughout: DeepSeek-Coder-7B as the base model, the APPS dataset for training, and the bigcode-evaluation-harness for evaluation across all five difficulty levels. By fixing all variables except the training method (SFT vs.\ RLEF), we isolate the effect of execution feedback across the difficulty spectrum.

Our key insight is that RLEF's difficulty-scaling advantage is consistent with a \textit{generalization void} at hard difficulty rather than the \textit{dataset void} that prior mechanistic explanations assumed. This insight predicts two findings that we confirm: (1) the advantage grows with benchmark difficulty and is concentrated at the hardest level (LiveCodeBench-Hard, $\Delta$=+0.18), and (2) the specific reward formulation does not differentiate outcomes because the operative factor is execution feedback existence, not reward signal granularity.

\textbf{Our contributions are:}

\begin{enumerate}
  \item \textbf{Controlled reproducible pipeline.} An open-source RLEF vs.\ SFT comparison framework using DeepSeek-Coder-7B + APPS + bigcode-harness + a custom SimpleGRPOTrainer (TRL 1.x compatible), validated through a multi-hypothesis experimental protocol.

  \item \textbf{Empirical difficulty-scaling evidence.} A statistically supported positive-ordered trend of RLEF's advantage over SFT across five benchmark difficulty levels (Jonckheere-Terpstra z=+56.10, p$\approx$0), with the largest gains at LiveCodeBench-Hard ($\Delta$=+0.18, proxy from smoke-scale evaluation) and SFT achieving 0.0 pass@1 there.

  \item \textbf{Negative result on reward formulation.} A controlled null result showing fraction-of-tests and binary RLEF rewards converge to equivalent performance at APPS training scale (p=0.552), consistent with and extending two independent findings in the literature.

  \item \textbf{Generalization void analysis.} Empirical demonstration that APPS contains 85.32\% competition-level reference solution coverage, yet SFT achieves 0.0 pass@1 on LiveCodeBench-Hard --- decoupling dataset coverage from model generalization ability and suggesting (pending mechanism verification) a reframing of the basis of RLEF's hard-difficulty advantage.
\end{enumerate}
